% TOP_PROBLEM: {"id": 2306066, "title": "Research Problems in Function Theory — Problem 6.66", "category_id": 8, "category": {"id": 8, "name": "analysis", "display_name": "Analysis", "description": "Limits, continuity, calculus, and function theory.", "slug": "analysis", "order_index": 8, "created_at": "2026-07-31T15:26:25.671Z"}, "status": "open"}
% FIRST_POSTED: null
% ENTRY_KIND: statement_only
% Imported from: batch03.tex; UnsolvedMath ID 2306066
\documentclass[11pt]{article}
\usepackage[margin=25mm]{geometry}
\usepackage{fontspec}
\setmainfont{Latin Modern Roman}
\usepackage{amsmath,amssymb,mathtools}
\usepackage{unicode-math}
\setmathfont{Latin Modern Math}
\usepackage{xurl,hyperref}
\hypersetup{colorlinks=true,urlcolor=blue}
\setcounter{secnumdepth}{0}
\setlength{\parindent}{0pt}
\setlength{\parskip}{5pt}
\setlength{\emergencystretch}{3em}
\newcommand{\sref}[2]{\hyperlink{source:#1:#2}{[S#2]}}
\newcommand{\cataloguescope}[1]{\paragraph{Recorded target.}#1}
\begin{document}
% UnsolvedMath ID 2306066; source catalogue code AMR-022-6066
\setcounter{section}{2306065}
\section{Research Problems in Function Theory — Problem 6.66}\label{2306066}
{\small\sffamily UnsolvedMath ID 2306066 · Analysis}
\subsection{Definitions and mathematical statement}

\paragraph{Shared notation.}$\mathbb N=\{1,2,\ldots\}$, and $\mathbb Z,\mathbb Q,\mathbb R,\mathbb C$ denote the integers, rationals, reals and complex numbers. Any other conventions are specified locally.


The exterior class $\Sigma$ consists of holomorphic injective functions on $\{z:|z|>1\}$ with Laurent expansion $f(z)=z+b_0+\sum_{j=1}^{\infty}b_jz^{-j}$. A function $f$ in a class $\mathcal A$ is an extreme point of that class if $f=tg+(1-t)h$, with $g,h\in\mathcal A$ and $0<t<1$, forces $g=h=f$. Let $\Sigma_0=\{g\in\Sigma:b_0=0\}$, and write $K_g=\mathbb C\setminus g(\{|z|>1\})$ for the compact omitted set. Measure means two-dimensional Lebesgue measure. This class-relative extreme-point convention and the distinct closed-convex-hull convention are stated separately in \sref{2306066}{3}.
\paragraph{Statement recorded in the supplied source.}
Characterize all extreme points of $\Sigma_0$. Is $g\in\Sigma_0$ an extreme point if and only if the omitted set $K_g$ has area zero? The source records the zero-area condition as sufficient; the question asks whether it is also necessary. \sref{2306066}{2}
Update 6.66 attributes an affirmative answer to Hamilton. \sref{2306066}{2}
\subsection{Short English statement}
Characterize extreme normalized exterior maps: must their omitted sets have zero area?
\subsection{Sources}
\begin{description}
\item[\hypertarget{source:2306066:1}{S1}] ULAM.AI and the UnsolvedMath Contributors, \emph{UnsolvedMath: A Curated Collection of Open Mathematics Problems}, record 2306066 (AMR-022-6066). CC BY 4.0. \url{https://huggingface.co/datasets/ulamai/UnsolvedMath}.
\item[\hypertarget{source:2306066:2}{S2}] W. K. Hayman and E. F. Lingham, \emph{Research Problems in Function Theory} (2018), Chapter 6, Problem 6.66 and its complete update; Chapter 6 initial notation. \url{https://arxiv.org/abs/1809.07200}.
\item[\hypertarget{source:2306066:3}{S3}] V. Allu and A. Pandey, \emph{Support points of some classes of analytic and univalent functions} (2022), Introduction and Preliminaries, p.3. \url{https://arxiv.org/pdf/2209.10607}.
\end{description}
\label{end:2306066}
\end{document}
